\documentclass[12pt,a4paper]{article}
\usepackage[utf8]{inputenc}
\usepackage[ngerman]{babel}
\usepackage{xcolor}
\usepackage{setspace}
\usepackage[margin=2.5cm]{geometry}

% Farbdefinitionen für grammatische Kategorien
\definecolor{verbcolor}{RGB}{0,128,0}      % Grün für Verbstämme
\definecolor{nouncolor}{RGB}{255,140,0}    % Orange für Nomen
\definecolor{conjcolor}{RGB}{255,0,0}      % Rot für Konjunktionen
\definecolor{prepcolor}{RGB}{0,0,255}      % Blau für Präpositionen/Partikel
\definecolor{annotcolor}{RGB}{128,128,128} % Grau für Annotationen
\definecolor{pluralcolor}{RGB}{255,0,255}  % Magenta für Plural
\definecolor{locationcolor}{RGB}{200,150,255} % Lila für Ortsangaben

% Makros für Verben und Verbformen
\newcommand{\verbstamm}[1]{\textcolor{verbcolor}{#1}}
\newcommand{\verbpart}[1]{\textcolor{prepcolor}{#1}}

% Makros für Nomen
\newcommand{\nomen}[1]{\textcolor{nouncolor}{#1}}

% Makros für Konjunktionen und Partikeln
\newcommand{\konj}[1]{\textcolor{conjcolor}{#1}}

% Makros für Präpositionen und Partikeln
\newcommand{\prep}[1]{\textcolor{prepcolor}{#1}}
\newcommand{\partikel}[1]{\textcolor{prepcolor}{#1}}

% Makros für besondere Formen
\newcommand{\plural}[1]{\textcolor{pluralcolor}{#1}}
\newcommand{\ort}[1]{\textcolor{locationcolor}{#1}}

% Makro für Annotationen in eckigen Klammern
\newcommand{\annot}[1]{{\footnotesize\textit{[#1]}}}

% Absatzformatierung
\setlength{\parindent}{0pt}
\setlength{\parskip}{1em}

\begin{document}

\section*{Textstruktur}

\subsection*{Josua 3,1-5}

\textbf{1} \konj{Da} \verbstamm{machte} \annot{Präteritum-Indikativ-3P-Aktiv-Singular} \nomen{Josua} \partikel{früh} \verbpart{auf} \annot{Verbpartikel}, ¶

\hspace*{1cm} \konj{und} \plural{sie} \annot{das-Volk} \verbstamm{zogen} \annot{Präteritum-Indikativ-3P-Plural} \prep{aus} \annot{Verbpartikel} \plural{Sittim}, ¶

\hspace*{1cm} \konj{und} \verbstamm{kamen} \annot{Präteritum-Indikativ, Aktiv, 3P-Plural} \prep{an} \prep{den} \nomen{Jordan}, ¶

\hspace*{2cm} er \annot{Josua} \konj{und} \plural{alle} \nomen{Kinder} \nomen{Israels} \annot{das-Volk}; ¶

\hspace*{2cm} \konj{und} \plural{sie} \annot{Volk-und-Josua} \verbstamm{rasteten} \annot{Präteritum, Aktiv, Indikativ, 3p, Plural} \ort{dort}, \konj{ehe} \plural{sie} \annot{Volk-und-Josua} \verbstamm{hinüberzogen}. ¶

¶

\textbf{2} \konj{Nach} \prep{drei} \nomen{Tage} \konj{aber} ¶

\hspace*{1cm} \verbstamm{gingen} \prep{die} \nomen{Vorsteher} \prep{durch} \prep{das} \ort{Lager} ¶

\hspace*{2cm} \textbf{3} \konj{und} \verbstamm{geboten} \prep{dem} \nomen{Volk}: ¶

\hspace*{2cm} \konj{und} \verbstamm{sprachen}: ¶

\hspace*{3cm} \textbf{Wenn} \prep{ihr} \prep{die} \nomen{Bundeslade} \prep{des} \nomen{Herrn}, \partikel{eures} \nomen{Gottes}, \verbstamm{sehen} \verbstamm{werdet} ¶

\hspace*{4cm} \konj{und} \prep{die} \nomen{Priester}, \prep{die} \nomen{Leviten}, \prep{die} \plural{sie} \verbstamm{tragen}, ¶

\hspace*{4cm} \konj{so} \verbstamm{brecht} \verbpart{auf} \prep{von} \partikel{eurem} \nomen{Ort} ¶

\hspace*{5cm} \konj{und} \nomen{folgt} \prep{ihr} \verbpart{nach}! ¶

\hspace*{3cm} \textbf{4} \konj{Doch} \konj{soll} \prep{zwischen} \partikel{euch} ¶

\hspace*{3cm} \konj{und} \prep{ihr} \partikel{etwa} \nomen{2000} \nomen{Ellen} \nomen{Abstand} \verbstamm{sein}. ¶

\hspace*{2cm} \nomen{Kommt} \prep{ihr} \konj{nicht} \prep{zu} \ort{nahe}, ¶

\hspace*{2cm} \konj{damit} \prep{ihr} \prep{den} \nomen{Weg} \verbstamm{erkennt}, \prep{den} \prep{ihr} \verbstamm{gehen} \konj{sollt}; ¶

\hspace*{3cm} \konj{denn} \prep{ihr} \verbstamm{seid} \prep{den} \nomen{Weg} \konj{zuvor} \verbstamm{gegangen}! ¶

¶

\textbf{5} \konj{Und} \nomen{Josua} \verbstamm{sprach} \prep{zum} \nomen{Volk}: \nomen{Heiligt} \partikel{euch}, ¶

\vspace{2em}

\subsection*{Legende}

\begin{itemize}
    \item \verbstamm{Grün} -- Verbstämme und Verbformen
    \item \nomen{Orange} -- Nomen und Substantive
    \item \konj{Rot} -- Konjunktionen und Bindewörter
    \item \prep{Blau} -- Präpositionen, Artikel und Partikeln
    \item \plural{Magenta} -- Pluralformen und Pronomen
    \item \ort{Lila} -- Ortsangaben
    \item \annot{Grau, kursiv} -- Grammatische Annotationen
\end{itemize}

\vspace{1em}
\textit{¶ = Textsegmentierung/Strukturmarkierung}

\end{document}
